\section{Conclusion}\label{sec:conclusion}

We asked a budgeting question about adversarial trajectory planning: with
revisions rationed to a public calendar, how much budget is needed, where
should it be spent, and can a computation certify the answer? The paper's
answer has three parts.

The frontier is computable. For a history-dependent sealed-block game, the
value attainable with at most $K$ revisions can be computed exactly up to a
certified interval, and the frontier is monotone by policy inclusion. On the
frozen design the budget is not slack: the smallest budget attaining the
full-flexibility value is the full budget in most cells, so the binding
constraint is the number of revisions rather than the particular schedule.
Where the budget does bite less, the maximizing calendars concentrate on the
early engaged epochs --- an empirical property, not a consequence of a greedy
argument, because the calendar objective is not submodular.

The certificates are valid and their sharpness is reported. A constructive
deletion certificate produces a feasible coarse policy together with an upper
bound on its loss, and an interval-loss bound certifies a budget under a fully
flexible opponent by a shortest path over public intervals. Both are bounds,
not optimality claims: the interval bound maximises over histories and the
deletion bound concerns one constructed policy, so we report where each is
loose instead of presenting a bound as a tightness result.

The scope is bounded on purpose. The physical model is a leader--follower
line-ahead formation with a frozen directional kernel; the design is a fixed
computational benchmark rather than a statistical sample; and an earlier
hypothesis that the range-advantaged side values revisions differently was
falsified at this fidelity and is kept as a negative result. The contribution
is the certified frontier and the structure that supports it, not a new
general equilibrium-computation method.

The natural next questions are the ones the certificates cannot yet answer:
tight rather than merely valid bounds, a game in which both sides select
calendars, and a formation model that deforms.
